\chapter{Hausdorff 拓扑空间上的测度}

若 T 是一个集合，A 是 T 的一个子集，则只要不致引起混淆，我们就用 \(\varphi_{A}\) 表示 A 的特征函数。集合 \(\overline{R}_{+}^{T}\) ——T 上定义的（有限或非有限的）数值函数 \(\geqslant0\) 全体——将记作 \(\mathcal{F}_{+}(T)\) ，当关于 T 没有歧义时也简记为 \(F_{+}\) ；该集合总赋有其自然的序结构。回忆 \(F_{+}\) 中两个元素的乘积总有定义，这要归功于约定 \(0\cdot(+\infty)=0\) 。若 A 是 T 的子集，f 是定义在 T 上的函数，则只要不引起混淆，f 在 A 上的限制 \(f|_{A}\) 在本章中可记作 \(f_{A}\) ；对诱导测度将采用类似的记号。另一方面，若 \(f\in\mathcal{F}_{+}(A)\) ，我们将用 \(f^{0}\) 表示 f 到 T 的 0 延拓，即在 A 上与 f 相合、在 T-A 上与 0 相合的 T 上的函数。

本章所考虑的拓扑空间，除明确说明相反情形外，都假定是 Hausdorff 的。从 §1 的 No.~4 起，除 §5 外，所有测度除明确说明相反情形外都假定为正测度。

